% ---
% title: What it was like writing my thesis
% date: 2026-08
% category:
% nodes: work 
% links: three-terms-in-cambridge
% ---
\documentclass{article}
\usepackage[utf8]{inputenc}
\usepackage{amsmath, amssymb, amsthm}
\usepackage{geometry}
\usepackage{graphicx}
\usepackage{hyperref}

\geometry{a4paper, margin=1in}

\title{What is was like writing my thesis}
\date{August 2026}
\author{Daksh Mehta}

\begin{document}

IN PROGRESS

This week I submitted my thesis. 

I was first introduced to the concept of fractal writing by my research supervisor at King's College London, (add footnote here). Fractals are repeating patterns that unfurl with increasing complexity (fix). Applying this to writing looks something like this. 

Fully worked example

Talk through how each word in the starting sentence can be unpacked. 

Preferable to writing in sections like introduction, methods, results, discussion. Helps to keep the big picture in the centre of the writing. 

Also I find that it helps me with blocks. It is kind of like increased the surface area of your substrate. Instead of writing and trying to generate one sentence after another, it gives lots of little starting points that can help you make more consistent progress. If inspiration or ideas on a certain topic run out, can just pick up somewhere else. 

Also another thing that I did was try and focus on putting words on the page. I kept my fingers on the keyboard and typed. If I didn’t know exactly how a sentence should end, I wrote something approximate. If I knew that a reference was needed, I wrote (add ref). If something needed explaining, I wrote (explain this). If I knew that a paragraph was incomplete, I left a note and moved on.




\end{document}